\documentclass[11pt,a4paper]{article}
\usepackage{amsmath}
\usepackage{amssymb}
\usepackage{geometry}
\usepackage{graphicx}
\usepackage{hyperref}
\usepackage{booktabs}
\usepackage{bm}

\geometry{margin=1in}

\title{Vehicle Dynamics with Dugoff Tire Model:\\
From Kinematic to Dynamic Simulation}
\author{Mario Tilocca}
\date{\today}

\begin{document}

\maketitle

\begin{abstract}
This document presents the transition from a kinematic bicycle model to a dynamic vehicle model incorporating the Dugoff tire model. We derive the complete force-based dynamics, normal load distribution with load transfer, slip ratio computation, and friction-limited tire forces. This enables realistic simulation of traction limits, wheel slip, and combined longitudinal-lateral dynamics for autonomous mining haul truck applications.
\end{abstract}

\tableofcontents
\clearpage

\section{Introduction}

\subsection{Motivation for Dynamic Model}

The kinematic bicycle model assumes:
\begin{itemize}
    \item \textbf{No slip:} Pure rolling condition at all times
    \item \textbf{Unlimited traction:} Tire forces can be arbitrarily large
    \item \textbf{Instantaneous response:} No force buildup dynamics
\end{itemize}

These assumptions break down when:
\begin{itemize}
    \item Vehicle operates on low-friction surfaces (gravel, wet haul roads)
    \item Emergency braking or aggressive acceleration
    \item Combined maneuvers (braking while turning)
    \item Realistic traction control and ABS simulation required
\end{itemize}

The Dugoff tire model addresses these by:
\begin{itemize}
    \item Computing friction-limited forces based on normal load
    \item Modeling slip ratios (longitudinal and lateral)
    \item Capturing smooth saturation from linear to sliding regimes
\end{itemize}

\subsection{Model Architecture Evolution}

\textbf{Kinematic Model (V1):}
\begin{equation}
\text{Actuator} \to \text{Drive} \to \text{Velocity} \to \text{Position/Yaw}
\end{equation}

\textbf{Dynamic Model (V2):}
\begin{equation}
\text{Actuator} \to \text{Drive} \to \boxed{\text{Tire Forces}} \to \text{Acceleration} \to \text{Velocity} \to \text{Position/Yaw}
\end{equation}

where Tire Forces are computed by the Dugoff model based on slip and normal load.

\section{Force-Based Longitudinal Dynamics}

\subsection{Equations of Motion}

Newton's second law in the longitudinal direction:
\begin{equation}
m \frac{dv}{dt} = \sum F_x - F_{\text{drag}} - F_{\text{roll}}
\end{equation}

where:
\begin{align}
\sum F_x &= F_{x,\text{fl}} + F_{x,\text{fr}} + F_{x,\text{rl}} + F_{x,\text{rr}} \quad \text{(sum of tire forces)} \\
F_{\text{drag}} &= \frac{1}{2} \rho C_d A v^2 \\
F_{\text{roll}} &= C_{\text{roll}}
\end{align}

\subsection{Tire Force Allocation}

For a rear-wheel-drive vehicle:
\begin{align}
F_{x,\text{fl}} &= 0 \quad \text{(front wheels: unpowered)} \\
F_{x,\text{fr}} &= 0 \\
F_{x,\text{rl}} &= f_{\text{Dugoff}}(\sigma_{x,\text{rl}}, \sigma_{y,\text{rl}}, F_{z,\text{rl}}) \\
F_{x,\text{rr}} &= f_{\text{Dugoff}}(\sigma_{x,\text{rr}}, \sigma_{y,\text{rr}}, F_{z,\text{rr}})
\end{align}

where $f_{\text{Dugoff}}$ is the tire model function (detailed in Section 4).

\subsection{Comparison: Kinematic vs Dynamic}

\begin{table}[h]
\centering
\begin{tabular}{lcc}
\toprule
\textbf{Property} & \textbf{Kinematic} & \textbf{Dynamic} \\
\midrule
Traction limit & Unlimited & $\mu F_z$ \\
Wheel slip & None & Computed \\
Braking distance & Underestimated & Realistic \\
Surface friction & Not modeled & $\mu \in [0.1, 0.9]$ \\
Computational cost & Low & Medium \\
\bottomrule
\end{tabular}
\caption{Kinematic vs dynamic model comparison}
\end{table}

\section{Normal Load Distribution}

\subsection{Static Load Distribution}

Weight distribution between front and rear axles:
\begin{align}
F_{z,\text{front}}^{\text{static}} &= \frac{mg \cdot d_{\text{rear}}}{L} \\
F_{z,\text{rear}}^{\text{static}} &= \frac{mg \cdot d_{\text{front}}}{L}
\end{align}

where:
\begin{itemize}
    \item $m$ -- vehicle mass (kg)
    \item $g = 9.81$ m/s$^2$ -- gravitational acceleration
    \item $L$ -- wheelbase (m)
    \item $d_{\text{front}}$ -- distance from CoG to front axle (m)
    \item $d_{\text{rear}}$ -- distance from CoG to rear axle (m)
\end{itemize}

Constraint: $d_{\text{front}} + d_{\text{rear}} = L$.

\subsection{Longitudinal Load Transfer}

During acceleration or braking, normal load shifts between axles:
\begin{equation}
\Delta F_z = \frac{m a_x h_{\text{CoG}}}{L}
\end{equation}

where:
\begin{itemize}
    \item $a_x$ -- longitudinal acceleration (m/s$^2$), positive = acceleration
    \item $h_{\text{CoG}}$ -- center of gravity height (m)
\end{itemize}

\textbf{Dynamic normal loads:}
\begin{align}
F_{z,\text{front}} &= F_{z,\text{front}}^{\text{static}} - \Delta F_z \\
F_{z,\text{rear}} &= F_{z,\text{rear}}^{\text{static}} + \Delta F_z
\end{align}

\textbf{Physical interpretation:}
\begin{itemize}
    \item $a_x > 0$ (acceleration): Weight transfers to rear, $F_{z,\text{rear}}$ increases
    \item $a_x < 0$ (braking): Weight transfers to front, $F_{z,\text{front}}$ increases
\end{itemize}

\subsection{Per-Wheel Normal Loads}

Assuming equal left/right distribution (no lateral load transfer in V2):
\begin{align}
F_{z,\text{fl}} &= F_{z,\text{fr}} = \frac{F_{z,\text{front}}}{2} \\
F_{z,\text{rl}} &= F_{z,\text{rr}} = \frac{F_{z,\text{rear}}}{2}
\end{align}

\textbf{Note:} Lateral load transfer during cornering is a future enhancement (V3).

\subsection{Example: XCMG XDE320 Mining Truck}

\textbf{Parameters:}
\begin{align}
m &= 218{,}000 \text{ kg (empty)} \\
L &= 6.3 \text{ m} \\
h_{\text{CoG}} &= 3.2 \text{ m} \\
d_{\text{front}} &= 2.52 \text{ m (40\% of wheelbase)} \\
d_{\text{rear}} &= 3.78 \text{ m (60\% of wheelbase)}
\end{align}

\textbf{Static loads:}
\begin{align}
F_{z,\text{front}}^{\text{static}} &= \frac{218{,}000 \times 9.81 \times 3.78}{6.3} = 1{,}288{,}000 \text{ N} \\
F_{z,\text{rear}}^{\text{static}} &= \frac{218{,}000 \times 9.81 \times 2.52}{6.3} = 862{,}000 \text{ N}
\end{align}

\textbf{Load transfer during braking at $a_x = -2.0$ m/s$^2$:}
\begin{align}
\Delta F_z &= \frac{218{,}000 \times (-2.0) \times 3.2}{6.3} = -222{,}000 \text{ N} \\
F_{z,\text{front}} &= 1{,}288{,}000 - (-222{,}000) = 1{,}510{,}000 \text{ N} \\
F_{z,\text{rear}} &= 862{,}000 + (-222{,}000) = 640{,}000 \text{ N}
\end{align}

Result: Front axle load increases by 17\%, rear decreases by 26\% during braking.

\section{Dugoff Tire Model}

\subsection{Fundamental Equations}

Tire forces are given by:
\begin{align}
F_x &= C_x \sigma_x f(\lambda) \label{eq:Fx_dugoff}\\
F_y &= C_y \sigma_y f(\lambda) \label{eq:Fy_dugoff}
\end{align}

where:
\begin{itemize}
    \item $F_x, F_y$ -- longitudinal and lateral tire forces (N)
    \item $C_x, C_y$ -- slip stiffness coefficients (N), load-dependent
    \item $\sigma_x, \sigma_y$ -- longitudinal and lateral slip ratios
    \item $f(\lambda)$ -- friction saturation function
    \item $\lambda$ -- friction utilization parameter
\end{itemize}

\subsection{Slip Ratios}

\textbf{Longitudinal slip ratio:}
\begin{equation}
\sigma_x = \frac{\omega R - V_x}{V_x}
\end{equation}

where:
\begin{itemize}
    \item $\omega$ -- wheel angular velocity (rad/s)
    \item $R$ -- tire effective radius (m)
    \item $V_x$ -- vehicle longitudinal velocity (m/s)
\end{itemize}

\textbf{Sign convention:}
\begin{align}
\sigma_x > 0 &\quad \text{Driving (wheel faster than vehicle)} \\
\sigma_x < 0 &\quad \text{Braking (wheel slower than vehicle)} \\
\sigma_x = 0 &\quad \text{Free rolling}
\end{align}

\textbf{Lateral slip ratio:}
\begin{equation}
\sigma_y = \frac{V_y}{V_x} \approx \tan(\alpha) \approx \alpha
\end{equation}

where $\alpha$ is tire slip angle (rad), valid for small angles ($|\alpha| < 0.2$ rad).

\subsection{Load-Dependent Slip Stiffness}

Tire stiffness increases with normal load following an empirical power law:
\begin{align}
C_x(F_z) &= C_{x,\text{base}} \left(\frac{F_z}{F_{z,\text{ref}}}\right)^{n} \\
C_y(F_z) &= C_{y,\text{base}} \left(\frac{F_z}{F_{z,\text{ref}}}\right)^{n}
\end{align}

where:
\begin{itemize}
    \item $C_{x,\text{base}}, C_{y,\text{base}}$ -- baseline stiffness at reference load (N)
    \item $F_{z,\text{ref}}$ -- reference normal load (N), typically 800 kN for mining tires
    \item $n$ -- load exponent, typically 0.5 (square-root dependence)
\end{itemize}

\textbf{Rationale:} Contact patch area increases with load, leading to increased grip.

\subsection{Friction Utilization Parameter}

The key parameter governing saturation:
\begin{equation}
\lambda = \frac{\mu F_z}{2\sqrt{C_x^2 \sigma_x^2 + C_y^2 \sigma_y^2}}
\end{equation}

where $\mu$ is the coefficient of friction (surface-dependent).

\textbf{Physical interpretation:}
\begin{align}
\lambda \geq 1 &\quad \text{Linear regime (no saturation)} \\
\lambda < 1 &\quad \text{Nonlinear regime (friction-limited)}
\end{align}

\subsection{Friction Saturation Function}

\begin{equation}
f(\lambda) = \begin{cases}
(2 - \lambda)\lambda = 2\lambda - \lambda^2, & \lambda < 1 \\
1, & \lambda \geq 1
\end{cases}
\end{equation}

\textbf{Properties:}
\begin{itemize}
    \item Continuous at $\lambda = 1$: $\lim_{\lambda \to 1^-} f(\lambda) = 1$
    \item Smooth transition from linear to saturated regime
    \item Maximum value: $f_{\max} = 1$ at $\lambda = 1$
\end{itemize}

\subsection{Proof: Force Saturation}

\textbf{Theorem:} The total tire force never exceeds the friction limit:
\begin{equation}
\sqrt{F_x^2 + F_y^2} \leq \mu F_z
\end{equation}

\textbf{Proof:} For $\lambda < 1$ (friction-limited case):
\begin{align}
\sqrt{F_x^2 + F_y^2} &= \sqrt{(C_x \sigma_x f)^2 + (C_y \sigma_y f)^2} \\
&= f(\lambda) \sqrt{C_x^2 \sigma_x^2 + C_y^2 \sigma_y^2} \\
&= (2\lambda - \lambda^2) \cdot \frac{\mu F_z}{2\lambda} \quad \text{(from definition of $\lambda$)} \\
&= \left(2 - \lambda\right) \frac{\mu F_z}{2} \\
&< \mu F_z \quad \text{(since } 0 < \lambda < 1\text{)}
\end{align}

For $\lambda \geq 1$, the tire is not friction-limited by construction. $\square$

\section{Lateral Velocity Computation}

\subsection{Front Axle Lateral Velocity}

Lateral velocity at front tire contact patch:
\begin{equation}
V_{y,\text{front}} = V_y + \dot{\psi} \cdot d_{\text{front}} + V_x \sin(\delta)
\end{equation}

where:
\begin{itemize}
    \item $V_y$ -- vehicle lateral velocity at CoG (m/s)
    \item $\dot{\psi}$ -- yaw rate (rad/s)
    \item $\delta$ -- steering angle (rad)
\end{itemize}

For small angles: $\sin(\delta) \approx \delta$.

\subsection{Rear Axle Lateral Velocity}

\begin{equation}
V_{y,\text{rear}} = V_y - \dot{\psi} \cdot d_{\text{rear}}
\end{equation}

\subsection{Yaw Rate from Bicycle Model}

From kinematic bicycle model:
\begin{equation}
\dot{\psi} = \frac{V_x}{L} \tan(\delta) \approx \frac{V_x \delta}{L}
\end{equation}

\section{Implementation Algorithm}

\subsection{Computation Sequence}

\textbf{Step 1: Compute Normal Loads}
\begin{enumerate}
    \item Static distribution: $F_{z,\text{front}}^{\text{static}}, F_{z,\text{rear}}^{\text{static}}$
    \item Load transfer: $\Delta F_z = \frac{m a_x h_{\text{CoG}}}{L}$
    \item Dynamic loads: $F_{z,\text{front}} = F_{z,\text{front}}^{\text{static}} - \Delta F_z$
    \item Per-wheel: $F_{z,\text{fl}} = F_{z,\text{fr}} = F_{z,\text{front}} / 2$
\end{enumerate}

\textbf{Step 2: Compute Slip Ratios}
\begin{enumerate}
    \item Get wheel angular velocities: $\omega_{\text{fl}}, \omega_{\text{fr}}, \omega_{\text{rl}}, \omega_{\text{rr}}$
    \item Longitudinal slip: $\sigma_x = (\omega R - V_x) / V_x$
    \item Lateral velocity: $V_y = V_y(\dot{\psi}, \delta)$ from Section 5
    \item Lateral slip: $\sigma_y = V_y / V_x$
\end{enumerate}

\textbf{Step 3: Compute Tire Forces (Dugoff)}
\begin{enumerate}
    \item Load-dependent stiffness: $C_x(F_z), C_y(F_z)$
    \item Friction parameter: $\lambda = \mu F_z / (2\sqrt{C_x^2 \sigma_x^2 + C_y^2 \sigma_y^2})$
    \item Saturation function: $f(\lambda)$
    \item Forces: $F_x = C_x \sigma_x f(\lambda)$, $F_y = C_y \sigma_y f(\lambda)$
\end{enumerate}

\textbf{Step 4: Update Dynamics}
\begin{enumerate}
    \item Total longitudinal force: $\sum F_x = F_{x,\text{rl}} + F_{x,\text{rr}}$
    \item Net force: $F_{\text{net}} = \sum F_x - F_{\text{drag}} - F_{\text{roll}}$
    \item Acceleration: $a = F_{\text{net}} / m$
    \item Velocity: $v_{k+1} = v_k + a \Delta t$
\end{enumerate}

\subsection{Discrete-Time Update}

Complete state update for timestep $\Delta t$:
\begin{align}
a_k &= \frac{F_{x,\text{rl}} + F_{x,\text{rr}} - F_{\text{drag}} - F_{\text{roll}}}{m} \\
v_{k+1} &= v_k + \Delta t \cdot a_k \\
\dot{\psi}_k &= \frac{v_k}{L} \tan(\delta_k) \\
\psi_{k+1} &= \psi_k + \Delta t \cdot \dot{\psi}_k \\
x_{k+1} &= x_k + \Delta t \cdot v_k \cos(\psi_k) \\
y_{k+1} &= y_k + \Delta t \cdot v_k \sin(\psi_k)
\end{align}

\section{Surface Friction Coefficients}

\subsection{Mining Haul Road Conditions}

\begin{table}[h]
\centering
\begin{tabular}{lcc}
\toprule
\textbf{Surface} & $\mu_{\text{peak}}$ & $\mu_{\text{slide}}$ \\
\midrule
Dry pavement (smooth) & 0.85 & 0.75 \\
Gravel compact (baseline) & 0.72 & 0.65 \\
Gravel loose & 0.55 & 0.48 \\
Iron ore dust (dry) & 0.45 & 0.38 \\
Iron ore dust (wet) & 0.30 & 0.25 \\
Mud & 0.25 & 0.20 \\
Sand (loose) & 0.40 & 0.35 \\
\bottomrule
\end{tabular}
\caption{Friction coefficients for Pilbara mining conditions}
\end{table}

\subsection{Velocity-Dependent Friction (Optional)}

Friction reduces at higher speeds due to reduced tire-road contact time:
\begin{equation}
\mu_{\text{effective}}(V) = \mu_{\text{peak}} \cdot \left(1 - k_v |V|\right)
\end{equation}

where $k_v = 0.003$ s/m (typical fade factor).

Constraint: $\mu_{\text{effective}} \geq 0.70 \mu_{\text{peak}}$ (minimum 70\% retention).

\section{Validation and Benchmarks}

\subsection{Expected Behavior}

\textbf{Scenario 1: Acceleration on gravel ($\mu = 0.72$)}

XCMG XDE320 loaded (538 tons):
\begin{align}
F_{z,\text{rear}} &\approx 0.75 \times 538{,}000 \times 9.81 = 3{,}956{,}000 \text{ N} \\
F_{x,\text{max}} &= 0.72 \times 3{,}956{,}000 = 2{,}848{,}000 \text{ N} \\
F_{x,\text{motor}} &= \frac{145{,}000}{1.93} = 75{,}100 \text{ N}
\end{align}

\textbf{Result:} Motor force $\ll$ traction limit → power-limited, not traction-limited.

\textbf{Scenario 2: Emergency braking on wet dust ($\mu = 0.30$)}

\begin{align}
F_{x,\text{max}} &= 0.30 \times 538{,}000 \times 9.81 = 1{,}585{,}000 \text{ N} \\
F_{x,\text{brake}} &= \frac{180{,}000}{1.93} = 93{,}300 \text{ N}
\end{align}

\textbf{Result:} Still power-limited, but much closer to traction limit.

\subsection{Stopping Distance Comparison}

\textbf{Kinematic model:} Assumes perfect traction, underestimates distance.

\textbf{Dynamic model:} Friction-limited, realistic distance:
\begin{equation}
d_{\text{stop}} = \frac{v_0^2}{2 \mu g}
\end{equation}

For $v_0 = 64$ km/h = 17.78 m/s on gravel ($\mu = 0.72$):
\begin{equation}
d_{\text{stop}} = \frac{17.78^2}{2 \times 0.72 \times 9.81} = 22.4 \text{ m}
\end{equation}

\section{Future Enhancements}

\subsection{Short-Term (V2.5)}
\begin{itemize}
    \item Lateral load transfer during cornering
    \item Combined longitudinal-lateral friction circle
    \item Wheel dynamics (rotational inertia)
\end{itemize}

\subsection{Medium-Term (V3)}
\begin{itemize}
    \item Pacejka "Magic Formula" tire model (higher fidelity)
    \item Suspension dynamics (vertical motion, damping)
    \item Thermal tire effects
\end{itemize}

\subsection{Long-Term (V4)}
\begin{itemize}
    \item 3D terrain interaction (grades, banking)
    \item Tire wear modeling
    \item Multi-physics coupling (thermal-mechanical)
\end{itemize}

\section{Conclusion}

The Dugoff tire model provides a significant improvement over the kinematic bicycle model by:
\begin{itemize}
    \item Enforcing realistic traction limits based on surface friction
    \item Computing slip ratios and forces explicitly
    \item Enabling ABS, traction control, and stability control validation
    \item Supporting realistic braking distance and emergency maneuver simulation
\end{itemize}

This forms the foundation for dynamic vehicle simulation suitable for autonomous mining vehicle development.

\begin{thebibliography}{9}

\bibitem{dugoff1970}
Dugoff, H., Fancher, P. S., and Segel, L. (1970).
\textit{An Analysis of Tire Traction Properties and Their Influence on Vehicle Dynamic Performance}.
SAE Technical Paper 700377.

\bibitem{rajamani2012}
Rajamani, R. (2012). \textit{Vehicle Dynamics and Control}. Springer.

\bibitem{gillespie1992}
Gillespie, T. D. (1992). \textit{Fundamentals of Vehicle Dynamics}. SAE International.

\bibitem{pacejka2012}
Pacejka, H. (2012). \textit{Tire and Vehicle Dynamics}, 3rd Edition. Butterworth-Heinemann.

\end{thebibliography}

\end{document}